\section{Conclusion}\label{sec:conclusion}

We have presented the v1 mathematical scaffolding of BEAMFS, a
recovery calculus for filesystem resilience under Single-Event
Effects. The contribution of this report is not an implementation,
nor a measurement campaign: it is the formal apparatus that
subsequent versions will instantiate, measure, and prove against.

The calculus is constructed on a multi-layer physical threat model
covering DRAM, SRAM and NAND-flash, with explicit positioning of
emerging memory technologies. The five invariants (I1--I5) are
stated minimally and mapped directly to kernel constructs. The
soundness theorem (\cref{thm:soundness}) shows that the recovery
operator drives every consistent state, perturbed by any event
within the bounded perturbation family, either to a consistent
successor or to fail-closed in a finite number of iterations.
Silent corruption is excluded by construction.

The empirical content of v1 is necessarily indirect: a re-reading
of the FTRFS v1 dataset through the calculus, and a projected
recovery surface from canonical cross-section data. The next
versions of the report will close this gap. v2 will document the
FTRFS Phase~E fault-injection campaign and a corresponding M6
bench metric. v3 will explore proof-assistant mechanization of
\cref{thm:soundness} and a verified-compilation chain for the
operator. Beam-time at a recognized accelerator facility is the
measurement programme this report is intended to support, with
the present formal apparatus as scientific dossier.

This is also v1 in another sense: it establishes authorship and
fixes the scope. Subsequent versions will refine; the present
fixed point is the right place to draw a line and submit it to
public review.
